\documentclass[margin=10pt]{standalone}
\usepackage{tikz}
\usetikzlibrary{shapes.geometric, arrows}

\begin{document}

% Define block styles
\tikzstyle{block} = [draw, fill=blue!10, rectangle, 
    minimum height=3em, minimum width=4em]
\tikzstyle{sum} = [draw, fill=blue!10, circle, node distance=1cm]
\tikzstyle{input} = [coordinate]
\tikzstyle{output} = [coordinate]
\tikzstyle{pinstyle} = [pin edge={to-,thin,black}]

\begin{tikzpicture}[auto, node distance=2cm,>=latex']
    % --- Nodes ---
    % Start with the input reference
    \node [input, name=input] {};
    
    % Summing junction
    \node [sum, right of=input] (sum) {};
    
    % Controller Block (C(s))
    \node [block, right of=sum] (controller) {$C(s)$};
    
    % Plant Block (G(s))
    \node [block, right of=controller, node distance=3cm] (system) {$G(s)$};
    
    % Output node
    \node [output, right of=system, node distance=2cm] (output) {};
    
    % Feedback Block (H(s))
    \node [block, below of=system] (feedback) {$H(s)$};

    % --- Connections ---
    % Path from input to sum
    \draw [draw,->] (input) -- node {$r(t)$} (sum);
    
    % Path from sum to controller
    \draw [->] (sum) -- node {$e(t)$} (controller);
    
    % Path from controller to plant
    \draw [->] (controller) -- node {$u(t)$} (system);
    
    % Path from plant to output
    \draw [->] (system) -- node [name=y] {$y(t)$}(output);
    
    % Feedback branch point and connection to H(s)
    \draw [->] (y) |- (feedback);
    
    % Path from H(s) back to sum
    \draw [->] (feedback) -| node[pos=0.99] {$-$} node [near end] {$b(t)$} (sum);
    
\end{tikzpicture}

\end{document}
